\documentclass[10pt,a4paper]{amsart}
\usepackage[margin=19mm]{geometry}
\usepackage{amsmath,amssymb,mathtools}
\usepackage[scr=rsfs,cal=euler]{mathalfa}
\usepackage{xfrac}
\usepackage[unicode,hidelinks,bookmarks=false]{hyperref}
\newcommand{\ie}{i.e.\ }
\newcommand{\cf}{cf.\ }
\newcommand{\resp}{respectively\ }
\newcommand{\Subsection}{\subsection}
\providecommand{\nobreakdash}{\nobreak\hbox{-}}
\setlength{\emergencystretch}{2em}
\multlinegap0pt
\usepackage{bm}
\usepackage{bbm}
\usepackage{dsfont}
\usepackage{xspace}
\usepackage{cancel}
\usepackage{mathtools}
\usepackage{amsthm}
\makeatletter
\@ifundefined{theorem}{\newtheorem{theorem}{Theorem}}{}
\@ifundefined{lemma}{\newtheorem{lemma}[theorem]{Lemma}}{}
\@ifundefined{proposition}{\newtheorem{proposition}[theorem]{Proposition}}{}
\makeatother
\title[Volume-Independent Spectral Stability of Energy-Truncated Ef]{Volume-Independent Spectral Stability of Energy-Truncated Effective Hamiltonians in Quantum Spin Systems}
\author{Ayumi Ukai}
\date{Source: arXiv:2605.07410v1 (CC BY 4.0)}
\begin{document}
\maketitle

\noindent\emph{Source and licence.} Volume-Independent Spectral Stability of Energy-Truncated Effective Hamiltonians in Quantum Spin Systems. Version arXiv:2605.07410v1 (CC BY 4.0), distributed under the Creative Commons Attribution 4.0 International licence, as recorded in the arXiv licence metadata for this exact version and shown on the abstract page https://arxiv.org/abs/2605.07410v1. Full rights notice: licenses/arxiv-2605.07410-CC-BY-4.0.txt.\par\smallskip

\noindent\textit{Editorial scope.} Counted unit: the Lemma whose source label is \texttt{lem\_G1} in the article (source0.tex lines 1263--1271, source label \texttt{lem\_G1}), delivered together with the article's own complete original proof of it (source0.tex lines 1272--1343, one \texttt{proof} environment of 72 lines). No mathematics is written, repaired or re-derived: statement and proof are the article's own text line for line. Required context retained uncounted: thm\_main\_finite (lines 261--284); lem\_AKL2.1\_5.1 (lines 363--371); thm\_main\_infinite (lines 668--691); prop\_AKL (lines 746--777); lem\_derivation (lines 848--883); lem\_infinite2 (lines 1253--1258). Every definition, display and label of those lines is retained unchanged. Omitted from this package and not redistributed: 1--260, 285--362, 372--667, 692--745, 778--847, 884--1252, 1259--1262, 1344--1590. Nothing inside a retained excerpt is omitted or altered: every retained body is delivered exactly as the article's lines are written, so no omission receipt of any kind accompanies this package. Citations: the retained excerpts carry no citation command at all, so no bibliography entry of the article is transcribed and no reference is redistributed. Reference adaptation: none was needed and none was made; every reference inside the delivered unit resolves inside it, so no unresolved reference or citation is produced. Printed slips kept literal: no printed word, symbol, hypothesis or number of the counted statement or of its proof was altered. Layout and class: the article's own class is \texttt{article} with packages including geometry, fontenc, inputenc, lmodern, amsmath, amssymb, amsthm, mathtools, mathrsfs, xcolor, hyperref; the wrapper is the lane's pinned \texttt{amsart} editorial wrapper at 10pt on A4, which restates the theorem-like environments the retained text needs and adds the article's own macro definitions that the retained text needs (none), with extra wrapper packages (bm, bbm, dsfont, xspace, cancel, mathtools). The delivered numbering is the unit's own. Assets: no retained span contains a figure environment, a graphics-inclusion command, a PDF-page inclusion, a table or a TikZ picture; no image file is copied into this package, the package declares no asset dependency and no image file is redistributed with it. Licence: version arXiv:2605.07410v1 of this article is distributed under the Creative Commons Attribution 4.0 International licence, as recorded in the arXiv licence metadata for this exact version and shown on the abstract page https://arxiv.org/abs/2605.07410v1. A full rights notice is delivered at licenses/arxiv-2605.07410-CC-BY-4.0.txt. Characters: every retained excerpt is pure ASCII, so the delivered engine, XeLaTeX with its default Latin Modern text font, reports no missing character.
\begin{theorem}\label{thm_main_finite}
We set
\[
\lambda:=\frac{1}{4\mathfrak{j}_\Phi(N_\Phi+|\partial L|)},\quad N_\Phi:=\sup_{x\in\Gamma} |\{ y\in\Gamma \,;\, d(x,y)\leq \mathfrak{r}_\Phi\}|<\infty,
\]
\begin{align*}
\delta(p,q)&:=2\sqrt{2}(M+5\|H_{\partial L}\|+q) \exp(-\lambda(M-2p-18\|H_{\partial L}\|)),\\
\eta(\epsilon,\delta)&:=2\sqrt{2}(M+\|H_{\partial L}\|+\delta) \exp(-\lambda(M-2\epsilon-10\|H_{\partial L}\|)).
\end{align*}
\begin{itemize}
\item[(i)] For any $q>p\geq-2\|H_{\partial L}\|$, we have
\[
\left\|\Big(I-E^{H_\Lambda}(-\infty,q)\Big)\,
E^{\bar{H}_\Lambda}(-\infty,p]\right\|
\leq\frac{p+2\|H_{\partial L}\|+\delta(p,q)}{q+2\|H_{\partial L}\|}.
\]
\item[(ii)] For any $\delta>\epsilon\geq0$, we have
\[
\left\|\Big(I-E^{H_\Lambda}(-\delta,\delta)\Big)\,
E^{\bar{H}_\Lambda}[-\epsilon,\epsilon]\right\|
\leq\frac{\epsilon+\eta(\epsilon,\delta)}{\delta}.
\]
\end{itemize}
\end{theorem}

\begin{lemma}\label{lem_AKL2.1_5.1}
We set $\lambda$ as in Theorem \ref{thm_main_finite}. Then, for any $M,N\in\mathbb{R}$ and any $A\in\mathcal{A}_\Lambda$, we have
\begin{align*}
&\|E^{H_\Lambda}[M,\infty) \,A\, E^{H_\Lambda}(-\infty,N]\|\\
&\quad\leq \exp(-\lambda(M-N-4\|H_{\partial L}\|)) \left\|e^{\lambda(H'_L+H'_{L^c})}Ae^{-\lambda(H'_L+H'_{L^c})}\right\|,\\
&\|E^{H'_L+H'_{L^c}}[M,\infty) \,A\, E^{H'_L+H'_{L^c}}(-\infty,N]\|\\
&\quad\leq \exp(-\lambda(M-N-4\|H_{\partial L}\|)) \left\|e^{\lambda H_\Lambda}Ae^{-\lambda H_\Lambda}\right\|.
\end{align*}
\end{lemma}

\begin{theorem}\label{thm_main_infinite}
We set
\[
\lambda:=\frac{1}{4\mathfrak{j}_\Phi(N_\Phi+|X|)}, \quad N_\Phi:=\sup_{x\in\Gamma} |\{ y\in\Gamma \,;\, d(x,y)\leq \mathfrak{r}_\Phi\}|<\infty,
\]
\begin{align*}
\delta(p,q)&:=2\sqrt{2}(M+5\|H_X\|+q) \exp(-\lambda(M-2p-18\|H_X\|)),\\
\eta(\epsilon,\delta)&:=2\sqrt{2}(M+\|H_X\|+\delta) \exp(-\lambda(M-2\epsilon-10\|H_X\|)).
\end{align*}
\begin{itemize}
\item[(i)] For any $q>p\geq-2\|H_X\|$, we have
\[
\left\|\left(I-E^{H_\omega}(-\infty,q)\right)\,
E^{\bar{H}_\omega}(-\infty,p]\right\|
\leq\frac{p+2\|H_X\|+\delta(p,q)}{q+2\|H_X\|}.
\]
\item[(ii)] For any $\delta>\epsilon\geq0$, we have
\[
\left\|\left(I-E^{H_\omega}(-\delta,\delta)\right)\,
E^{\bar{H}_\omega}[-\epsilon,\epsilon]\right\|
\leq\frac{\epsilon+\eta(\epsilon,\delta)}{\delta}.
\]
\end{itemize}
\end{theorem}

\begin{proposition}[Infinite extension of Lemma \ref{lem_AKL2.1_5.1}]\label{prop_AKL}
We set $\lambda$ as in Theorem \ref{thm_main_infinite}. We assume that $A,B\in\mathcal{B}(\mathcal{H}_\omega)$ are such that, for every $s\in\mathbb{C}$, the sesquilinear forms
\[
(\phi,\psi)\in\mathcal{D}_{H_{X^c}}\times\mathcal{D}_{H_{X^c}}
\mapsto
\left\langle e^{\bar{s} H_{X^c}}\phi,
A e^{-sH_{X^c}}\psi\right\rangle,
\]
\[
(\phi,\psi)\in\mathcal D_{H_\omega}\times \mathcal D_{H_\omega}
\mapsto
\left\langle e^{\bar{s} H_\omega}\phi,
B e^{-sH_\omega} \psi\right\rangle
\]
extend to bounded sesquilinear forms on $\mathcal{H}_\omega\times\mathcal{H}_\omega$. We also assume that the corresponding bounded operators, denoted by
\[
e^{sH_{X^c}}Ae^{-sH_{X^c}},\quad e^{sH_\omega}Be^{-sH_\omega},
\]
define $\mathcal{B}(\mathcal H_\omega)$-valued entire functions of $s$.

Then for any $M,N\in\mathbb{R}$, we have
\begin{align*}
\|E^{H_\omega}[M,\infty) \,A\,
E^{H_\omega}(-\infty,N]\|
&\leq\exp(-\lambda(M-N-4\|H_X\|))\,
\|e^{\lambda H_{X^c}}Ae^{-\lambda H_{X^c}}\|,\\
\|E^{H_{X^c}}[M,\infty) \,B\,
E^{H_{X^c}}(-\infty,N]\|
&\leq\exp(-\lambda(M-N-4\|H_X\|))
\|e^{\lambda H_\omega}Be^{-\lambda H_\omega}\|.
\end{align*}
\end{proposition}

\begin{lemma}\label{lem_derivation}
We set $\lambda$ as in Theorem \ref{thm_main_infinite} and define
\[
S_X:=2\lambda,\quad R(S_X):=\{t\in\mathbb{C}\,;\,|t|<S_X\}.
\]
\begin{itemize}
\item[(i)] The series
\begin{align*}
\exp(s\delta_\Gamma)(H_X)
&:=\sum_{n=0}^\infty\frac{s^n}{n!} (\delta_\Gamma)^n(H_X),\\
\exp(s\delta_{X^c})(H_X)
&:=\sum_{n=0}^\infty\frac{s^n}{n!} (\delta_{X^c})^n(H_X)
\end{align*}
are absolutely convergent for any $s\in R(S_X)$. Moreover, we then have
\[
\max\{\|\exp(s\delta_\Gamma)(H_X)\|, \|\exp(s\delta_{X^c})(H_X)\|\}
\leq \frac{\|H_X\|}{1-|s|/S_X}.
\]

\item[(ii)] For any $s\in R(S_X)$ and any $A\in\mathcal{B}(\mathcal{H}_\omega)$, the following
$\mathcal{B}(\mathcal{H}_\omega)$-valued initial value problems posed on the interval $t\in[0,1]$ admit unique solutions:
\[
\frac{d}{dt} F_1(s,A,t)=[s\exp(ts\delta_\Gamma)(H_X),F_1(s,A,t)]
\quad\text{with}\quad F_1(s,A,0)=A,
\]
\[
\frac{d}{dt} F_2(s,A,t)=[s\exp(ts\delta_{X^c})(-H_X),F_2(s,A,t)]
\quad\text{with}\quad F_2(s,A,0)=A.
\]
Moreover, we then have
\[
\max\{\|F_1(s,A,t)\|,\|F_2(s,A,t)\|\}
\leq \exp\left(2|s|\frac{\|H_X\|}{1-t|s|/S_X}\right) \|A\|.
\]
\end{itemize}
\end{lemma}

\begin{lemma}\label{lem_infinite2}
We set $S_X$ as in Proposition \ref{prop_AKL}. Then, for any $s\in R(S_X)$ and any $\psi,\phi\in\mathcal{D}_{H_{X^c}}$, the following identity holds as a sesquilinear-form identity:
\[
\langle e^{\bar{s}H_{X^c}}\phi,H_X e^{-sH_{X^c}}\psi\rangle=\langle\phi,\exp(s\delta_{X^c})(H_X)\psi\rangle.
\]
\end{lemma}

\begin{lemma}\label{lem_G1}
For any $s\in R(S_X)$, we have
\[
\left\|
e^{s\bar{H}_{X^c}}H_X e^{-s\bar{H}_{X^c}}
\right\|
\leq 2\|H_X\| +\frac{\|H_X\|}{1-|s|/S_X}.
\]
\end{lemma}

\begin{proof}
First, we prove the statement for the case $\operatorname{Re}(s)\geq0$. For simplicity, we write
\[
E_1:=E^{H_{X^c}}(-\infty,-M],\quad
E_2:=E^{H_{X^c}}(-M,M),\quad
E_3:=E^{H_{X^c}}[M,\infty).
\]
Then, we decompose
\begin{align*}
e^{s\bar{H}_{X^c}}H_X e^{-s\bar{H}_{X^c}}
&= e^{s\bar{H}_{X^c}} H_X e^{-s\bar{H}_{X^c}}E_3
+ E_1 e^{s\bar{H}_{X^c}} H_X e^{-s\bar{H}_{X^c}}(E_1+E_2)\\
&\qquad+(E_2+E_3) e^{s\bar{H}_{X^c}} H_X e^{-s\bar{H}_{X^c}}(E_1+E_2).
\end{align*}

For the first and second terms, we have
\begin{align*}
&\max\left\{ \left\|e^{-s\bar{H}_{X^c}} E_3\right\|,\,
\left\|E_1 e^{s\bar{H}_{X^c}}\right\| \right\}\leq e^{-\operatorname{Re}(s) M},\\
&\max\left\{\left\| e^{s\bar{H}_{X^c}}\right\|,\,
\left\|e^{-s\bar{H}_{X^c}}\right\|\right\} \leq e^{\operatorname{Re}(s) M}. 
\end{align*}
This yields the bounds
\[
\max\left\{
\left\|
e^{s\bar{H}_{X^c}} H_X e^{-s\bar{H}_{X^c}} E_3
\right\|,
\left\|
E_1 e^{s\bar{H}_{X^c}} H_X e^{-s\bar{H}_{X^c}}
\right\|
\right\}
\leq \|H_X\|.
\]

For the third term, for any $\phi\in\mathcal{D}_{H_{X^c}}$, we have
\begin{align*}
\left\|
e^{-sH_{X^c}}e^{s\bar{H}_{X^c}}
(E_2+E_3)\phi
\right\|^2
&\leq \|E_2\phi\|^2+
\int_M^\infty e^{2\operatorname{Re}(s)(M-\mu)} \|E^{H_{X^c}}(d\mu)\phi\|^2\\
&\leq \|E_2\phi\|^2+ \|E_3\phi\|^2
=\|(E_2+E_3)\phi\|^2.
\end{align*}
Similarly, we have
\[
\left\|
e^{sH_{X^c}}e^{-s\bar{H}_{X^c}}(E_1+E_2)\phi
\right\|
\leq \|(E_1+E_2)\phi\|.
\]
By Lemma \ref{lem_infinite2} and Lemma \ref{lem_derivation}(i), we obtain
\begin{align*}
&\left\|
(E_2+E_3) e^{s\bar{H}_{X^c}} H_X e^{-s\bar{H}_{X^c}} (E_1+E_2)\phi
\right\|\\
&\leq \|\exp(s\delta_{X^c})(H_X)\| \|\phi\|
\leq \frac{\|H_X\|}{1-|s|/S_X} \|\phi\|.
\end{align*}
Since $\phi\in\mathcal{D}_{H_{X^c}}$ is arbitrary and $\mathcal{D}_{H_{X^c}}\subset\mathcal{H}_\omega$ is dense, this yields
\[
\left\|
(E_2+E_3) e^{s\bar{H}_{X^c}} H_X e^{-s\bar{H}_{X^c}} (E_1+E_2)
\right\|
\leq \frac{\|H_X\|}{1-|s|/S_X}.
\]
Putting all these together, we obtain the desired estimate for $\operatorname{Re}(s)\geq0$. 

The case $\operatorname{Re}(s)\leq0$ follows in the same way by exchanging $E_1$ and $E_3$. 
\end{proof}

\end{document}
